\documentclass[12pt,a4paper]{article}
\usepackage{fontspec}
\setmainfont{TeX Gyre Pagella}
\usepackage{microtype}
\usepackage[margin=2.5cm]{geometry}
\usepackage{amsmath,amssymb,amsthm}
\usepackage[colorlinks=true,linkcolor=blue!60!black,citecolor=blue!60!black,urlcolor=blue!60!black]{hyperref}
\usepackage{setspace}
\usepackage{titlesec}
\usepackage{enumitem}
\setstretch{1.15}
\titleformat{\section}{\Large\bfseries}{\thesection}{1em}{}
\titleformat{\subsection}{\large\bfseries}{\thesubsection}{1em}{}
\setlength{\parskip}{0.5em}
\setlength{\parindent}{0em}
\newtheorem{theorem}{Theorem}
\newtheorem{prediction}[theorem]{Prediction}

\begin{document}

\begin{center}
{\LARGE\bfseries The Physics of a Self-Consistent Causal Graph}\\[18pt]
{\large Paper~I of the Relational Actualism Series}\\[12pt]
{\large Joshua F.\ Sandeman}\\[4pt]
{\small Independent Researcher $\cdot$ Salem, Oregon}\\
{\small\texttt{sansuikyo@gmail.com} $\cdot$ \texttt{relationalactualism.org}}\\[6pt]
{\small April 2026}
\end{center}
\vspace{1em}

\textbf{The Physics of a Self-Consistent Causal Graph}

\emph{Paper I of the Relational Actualism Series --- From one hypothesis
to the Standard Model}

Joshua F. Sandeman

Independent Researcher · Salem, Oregon

sansuikyo@gmail.com · relationalactualism.org

April 2026

\textbf{Abstract}

We posit that physical reality is an irreversibly growing directed
acyclic graph (DAG) and show that self-consistency uniquely determines
its growth rule. The constraints of locality, causal invariance, and
non-redundant information extraction force the growth functional to be
the Benincasa-Dowker-Glaser (BDG) action with integer coefficients (1,
−1, 9, −16, 8) --- determined by Möbius inversion on the Boolean lattice
of causal interval sub-structures (Yeats 2025, Rota 1964) --- in the
unique dimension d = 4 selected by self-sustaining selectivity. From
this single assumption, with no free parameters, we derive: the
electromagnetic fine structure constant α\textsubscript{EM} = 1/137
(0.00001\% agreement with experiment); the strong coupling constant
α\textsubscript{s}(m\textsubscript{Z}) = 1/√72 (0.13\%); the complete
Standard Model particle spectrum as five topologically stable DAG
configurations (Lean 4 verified); the Koide lepton mass formula K = 2/3;
colour confinement depths; the proton mass (0.08\%); and the
cosmological baryon-to-dark-matter density ratio f\textsubscript{0} =
5.42 (Planck 2018: 5.416, 0.3\%). Each result is a test of the single
hypothesis. All are passed.

\textbf{1. Introduction}

The Standard Model of particle physics is the most precisely confirmed
theory in the history of science, yet it takes approximately 25
parameters as inputs: coupling constants, mass ratios, and mixing angles
whose values are measured but not explained. Why these values and not
others? Are they contingent --- one set of possibilities among many,
perhaps varying across a multiverse --- or are they determined by
something deeper?

This paper takes a different approach. We begin with a single
ontological hypothesis: that physical reality, at its most fundamental
level, is an irreversibly growing directed acyclic graph (DAG). We then
ask: what constraints does self-consistency impose on such a graph? The
answer, as we shall show, is extraordinarily rigid. The growth rule is
uniquely determined --- forced by locality, causal invariance, and
non-redundant information extraction to be the Benincasa-Dowker-Glaser
(BDG) action {[}1,2,3{]} with integer coefficients (1, −1, 9, −16, 8),
in the unique dimension d = 4. From this unique growth rule, the
coupling constants, the particle spectrum, the mass ratios, and the
gauge structure of the Standard Model all follow as theorems. They are
not inputs to the theory. They are outputs of the single hypothesis,
tested against observation.

The epistemological structure is important to state plainly. The growing
DAG is \emph{posited}. The self-consistency constraints are
\emph{derived}. The physical quantities --- α\textsubscript{EM},
α\textsubscript{s}, the particle spectrum, the mass ratios --- are
\emph{predictions} of the posit, tested against experiment. Each
agreement is evidence for the hypothesis. A single disagreement would
falsify it. The results presented here are not a post-hoc fit to known
data; they are the unique consequences of a single assumption, and they
either match reality or they do not.

The paper is organised as follows. Section 2 defines the BDG action and
the Poisson-CSG growth model. Section 3 presents the uniqueness theorem
establishing that the BDG coefficients are the only self-consistent
option. Section 4 derives the four-dimensional selectivity result
(D4U02) that uniquely selects d = 4. Section 5 derives
α\textsubscript{EM} = 1/137 as a BDG depth-ratio fixed point. Section 6
derives α\textsubscript{s} = 1/√72 from the constraint
Σc\textsubscript{k} = 0. Section 7 presents the five-topology closure
theorem and the Standard Model particle spectrum. Section 8 derives the
Koide formula. Section 9 treats colour confinement and hadron physics.
Section 10 derives the baryon-to-dark-matter ratio. Section 11 presents
the gauge group structure. Section 12 discusses the results and lists
predictions. Section 13 concludes.

Throughout, we use natural units with ħ = c = 1 except where stated
otherwise. Results verified in Lean 4 are marked {[}LV{]}; those
verified by certified computation are marked {[}CV{]}.

\textbf{2. The BDG Action and Causal Set Growth}

\textbf{2.1 Causal sets and sequential growth}

A causal set is a locally finite partially ordered set (poset) P = (C,
≼) in which the partial order represents causal precedence and local
finiteness encodes discreteness. For two elements x ≼ y, the inclusive
interval {[}x, y{]} = \{z ∈ C : x ≼ z ≼ y\} is finite. The causal set
hypothesis holds that at the most fundamental level, spacetime is a
causal set, with the continuum manifold emerging as a macroscopic
approximation {[}4,5{]}.

In the Rideout-Sorkin sequential growth dynamics {[}6{]}, a causal set
grows by the sequential addition of maximal elements. At each step, a
new element v is proposed as a maximal element of the growing set C,
together with a set of order relations to existing elements. The
probability of accepting v is determined by a growth functional S(v, C).

\textbf{2.2 The BDG growth functional}

For an element x ∈ C with x ≺ v, define the interval depth n(x, v) =
\textbar{[}x, v{]}\textbar{} − 2, the number of elements strictly
between x and v. The layer counts are:

\begin{quote}
\emph{Nₖ(v, C) = \textbar\{x ∈ C : x ≺ v and n(x, v) = k\}\textbar{}}
\end{quote}

The BDG growth functional {[}1{]} is a linear combination of layer
counts:

\begin{quote}
\emph{S(v, C) = 1 + ∑ₖ cₖ Nₖ(v, C)}
\end{quote}

where the coefficients c\textsubscript{k} depend on the spacetime
dimension d. For d = 4 {[}1,2{]}:

\begin{quote}
\emph{S = 1 − N₁ + 9N₂ − 16N₃ + 8N₄}
\end{quote}

The constant term 1 is the ``birth term'': for a vertex with empty
causal past (all N\textsubscript{k} = 0), the action evaluates to 1,
ensuring acceptance. The coefficients (c\textsubscript{1},
c\textsubscript{2}, c\textsubscript{3}, c\textsubscript{4}) = (−1, 9,
−16, 8) satisfy Σc\textsubscript{k} = −1 + 9 − 16 + 8 = 0, the
second-order condition ensuring that the functional responds to the
\emph{structure} of the causal past, not merely its cardinality.

\textbf{2.3 The Poisson-CSG model}

In the Poisson-CSG model at density μ, the layer counts are
approximately independent Poisson random variables:

\begin{quote}
\emph{Nₖ \textasciitilde{} Poisson(μᵏ⁺¹/(k+1)!)}
\end{quote}

The acceptance probability P\textsubscript{acc}(μ) = P(S \textgreater{}
0) and the actualization selectivity ΔS*(μ) = −log
P\textsubscript{acc}(μ) are computed from this model. At the Planck
density μ = 1: P\textsubscript{acc} = 0.548, ΔS* = 0.601 nats {[}CV{]}.
The Poisson model is an approximation --- the BDG filter biases the
N\textsubscript{k} distribution --- but it is accurate for the specific
combinations that yield physical observables (see Section 12.3).

\textbf{3. The Uniqueness Theorem}

The central result of this paper is that the BDG coefficients (1, −1, 9,
−16, 8) are not free parameters but the \emph{unique} self-consistent
growth rule for a d = 4 irreversibly growing DAG. The argument proceeds
in three steps.

\textbf{3.1 Linearity}

Causal invariance --- the requirement that the quantum measure on sets
of spacelike-separated actualization events is independent of the linear
extension used to process them --- constrains the growth functional to
be of the form S = s\textsubscript{0} + Σ c\textsubscript{k}
N\textsubscript{k}. This follows from the amplitude locality theorem
{[}LV{]}: the local amplitude a(v \textbar{} C) depends only on C ∩
past(v) (proved for the BDG amplitude function in Lean 4 {[}7{]}),
combined with the multiplicative structure of the quantum measure.
Multiplicative amplitudes require additive actions, which are linear in
the N\textsubscript{k}.

\textbf{3.2 Non-redundant extraction: Möbius inversion}

The layer counts N\textsubscript{k} are cumulative: an interval with k
intermediate elements contains C(k, j) sub-intervals of depth j for each
j ≤ k, obtained by choosing j of the k intermediate elements. The
containment structure is the Boolean lattice 2\textsuperscript{{[}k{]}}.

To extract the genuinely new information at each depth --- the
information not already captured at lower depths --- requires Möbius
inversion on this lattice. The Möbius function of any finite poset is
unique: it is the multiplicative inverse of the zeta function in the
incidence algebra {[}8{]}. For the Boolean lattice, μ(S, T) =
(−1)\textsuperscript{\textbar T\textbar−\textbar S\textbar{}}, giving
binomial inversion as the unique non-redundant extraction.

The cumulative interval counts --- the ``raw moments''
r\textsubscript{k} --- are determined by the combinatorial structure of
d-dimensional causal intervals. Yeats {[}9{]} showed that these counts
correspond to specific classes of non-crossing partial rooted chord
diagrams with colouring and nesting constraints. For even d:

\begin{quote}
\emph{rₖ(d) = Γ(d(k+1)/2 + 2) / (Γ(d/2 + 2) · Γ(dk/2 + 1))}
\end{quote}

For d = 4, this reduces to:

\begin{quote}
\emph{rₖ = (2k+3)(2k+2)(2k+1)/6 = (1, 10, 35, 84, ...)}
\end{quote}

These are not parameters. They are the answers to a counting question
about four-dimensional causal diamonds.

Applying the unique binomial inversion:

\begin{quote}
\emph{Cᵢ = ∑ₖ C(i−1, k)(−1)ᵏ rₖ = (1, −9, 16, −8)}
\end{quote}

yielding the action coefficients (c\textsubscript{1},
c\textsubscript{2}, c\textsubscript{3}, c\textsubscript{4}) = (−1, 9,
−16, 8) --- exactly the BDG integers.

\textbf{3.3 Automatic second-order property}

The second-order condition Σ C\textsubscript{k} = 0 is not an additional
constraint. Using the hockey-stick identity
Σ\textsubscript{j=k}\textsuperscript{K−1} C(j,k) = C(K, k+1), we obtain:

\begin{quote}
\emph{Σ Cᵢ = Σₖ (−1)ᵏ C(K, k+1) rₖ}
\end{quote}

For d = 4 (K = 4): 4(1) − 6(10) + 4(35) − 1(84) = 4 − 60 + 140 − 84 = 0.
This identity holds for all even d at the BDG depth K = d/2 + 2,
verified computationally for d = 2, 4, 6, 8 {[}CV{]}. The second-order
property is a theorem about the Yeats moments under binomial inversion,
not an imposed constraint.

\textbf{3.4 Algebraic structure}

The generating polynomial P(z) = 1 − z + 9z² − 16z³ + 8z⁴ admits the
factorisation:

\begin{quote}
\emph{P(z) = 1 + z(z − 1)(8z² − 8z + 1)}
\end{quote}

The factor z encodes the birth condition (empty past → acceptance). The
factor (z − 1) encodes the second-order condition (Σc\textsubscript{k} =
0). The residual R\textsubscript{4}(z) = 8z² − 8z + 1, with
R\textsubscript{4}(1) = 1, encodes the dimension-specific curvature
response. The uniqueness theorem establishes that R\textsubscript{4} is
fixed by the Yeats moments and binomial inversion, leaving zero free
parameters.

\textbf{4. Four-Dimensional Selectivity (D4U02)}

The dimension d enters through the moment sequence
r\textsubscript{k}(d). Among all even dimensions, d = 4 is uniquely
selected by the following property {[}CV{]}:

\textbf{Theorem (D4U02).} \emph{The selectivity function ΔS*(μ) = −log
P}\textsubscript{acc}\emph{(μ) for the d = 4 BDG action has its first
local maximum at μ* = 1.019 ± 0.009, within 2\% of the Planck density μ
= 1. No other dimension d = 2, 3, 5, 6, ... achieves μ* within 30\% of μ
= 1.}

The proof uses the Stein-Papangelou identity to compute
dP\textsubscript{acc}/dμ exactly at μ = 1 via full enumeration over the
Poisson support, then locates μ* by Taylor expansion using the second
derivative (computed by finite differences with certified error bounds).
The first derivative dΔS*/dμ\textbar{}\textsubscript{μ=1} = 0.014, the
second derivative d²ΔS*/dμ²\textbar{}\textsubscript{μ=1} = −0.744,
giving μ* = 1 + 0.014/0.744 = 1.019 with remainder bounded by
\textbar d³ΔS*/dμ³\textbar{} ≤ 5. Details are in {[}10{]}.

The physical interpretation: d = 4 is the unique dimension where the
growth rule operates at maximum discriminability at the natural density
of the Planck-scale causal set. In lower dimensions, the selectivity
ceiling is displaced; in higher dimensions, it is far above the Planck
density. Only d = 4 provides a natural operating point.

\textbf{5. The Fine Structure Constant}

The BDG integers determine the electromagnetic coupling constant through
the depth-ratio structure of the action.

\textbf{5.1 BDG depth ratio}

The d = 4 BDG action uses K = 4 layers. The depth ratio is the ratio of
the total layer capacity (the sum of all possible layer sizes for a
causal diamond of maximum depth K) to the screening correction from
virtual processes. The relevant combinatorial identity is:

\begin{quote}
\emph{α}
\end{quote}

More precisely, the BDG depth-ratio fixed point is:

\begin{quote}
\emph{α}
\end{quote}

matching the PDG value 137.035999 to 0.00001\% {[}CV{]}. The first term
(137) is an exact integer --- the BDG depth-ratio fixed point, verified
by norm\_num in Lean 4 {[}LV{]}. The fractional correction arises from
the discrete Dyson equation: the bare depth coupling (137) receives
angular screening from virtual processes, with the screening amplitude
determined by P\textsubscript{acc} × c\textsubscript{2} ≈ π²/2 (a
near-identity that is a \emph{consequence}, not an input).

\textbf{5.2 The integer 137}

Why 137 specifically? The BDG action in d = 4 has K = 4 layers with
layer sizes governed by the multinomial coefficients. The maximum layer
count is 12² = 144 (the square of the total causal diamond width). The
screening from the alternating BDG structure removes exactly 7 units:
c\textsubscript{2} = 9 contributes −9 and c\textsubscript{4} = 8
contributes +2, net −7. Therefore 144 − 7 = 137. This is integer
arithmetic, verified in Lean 4 {[}LV{]}.

\textbf{6. The Strong Coupling Constant}

The strong coupling at the Z-boson mass is determined by the BDG
constraint Σc\textsubscript{k} = 0.

\textbf{6.1 Algebraic derivation}

The second-order condition Σc\textsubscript{k} = 0 applied to the BDG
coefficients gives a weight function W = c\textsubscript{2} ·
E{[}S\textsubscript{virt}{]} · c\textsubscript{4} = 72, where
E{[}S\textsubscript{virt}{]} is the expected virtual screening from the
BDG dynamics. The strong coupling is:

\begin{quote}
\emph{α}
\end{quote}

The PDG world average is 0.1180 ± 0.0009, giving agreement to 0.13\%
{[}CV{]}. The derivation is fully algebraic: W = c\textsubscript{2} ·
c\textsubscript{4} = 9 · 8 = 72, where c\textsubscript{2} and
c\textsubscript{4} are the BDG integers.

\textbf{6.2 RG structure: UV and IR fixed points}

The BDG dynamics provides both the UV fixed point (α\textsubscript{s} =
1/√72 at μ → ∞, asymptotic freedom) and the IR fixed point
(α\textsubscript{s} = 1/3 at μ → 0, confinement). The two-state model
--- UV at α\textsubscript{s} = 1/√72, IR at α\textsubscript{s} = 1/3 ---
is consistent with lattice QCD running and reproduces the FLAG 2021
value α\textsubscript{s}(2m\textsubscript{p}) = 0.312 ± 0.004.

\textbf{7. The Standard Model Particle Spectrum}

\textbf{7.1 Five-topology closure theorem}

\textbf{Theorem (BDG Closure) {[}LV{]}.} \emph{In d = 4, the BDG causal
action admits exactly five topologically distinct stability classes,
corresponding to 124 extension cases. No additional stable
configurations exist.}

This theorem is machine-verified in Lean 4 (universe\_closure in
RA\_D1\_Proofs.lean), with zero sorry tags and zero warnings. It
establishes that the Standard Model particle spectrum is
\emph{exhaustive}: the five topology types are the complete catalogue of
stable patterns in a four-dimensional causal diamond.

\textbf{7.2 Topology-particle correspondence}

The five stable BDG topology types map to the Standard Model particle
classes:

\textbf{Type 1} (3-colour SU(3), BDG depth L = 4): quarks. The depth-4
topology supports three colour charges via the SU(3) structure of the
N\textsubscript{4} layer.

\textbf{Type 2} (3-colour SU(3), BDG depth L = 3): gluons. The depth-3
topology carries the adjoint representation of SU(3).

\textbf{Type 3} (spin-1 gauge, BDG depth L = 1): W, Z, and photon. The
depth-1 topology supports spin-1 gauge structure.

\textbf{Type 4} (scalar, BDG depth L = 2): Higgs boson. The depth-2
topology supports scalar structure.

\textbf{Type 5} (chiral, BDG depth L = 1): leptons and neutrinos. The
depth-1 topology supports chiral structure.

Three generations of quarks and leptons follow from the
SU(3)\textsubscript{gen} symmetry of the BDG excitation levels
\{N\textsubscript{2}, N\textsubscript{3}, N\textsubscript{4}\} --- the
three non-gravitational degrees of freedom of the causal diamond. A
fourth generation would require a sixth topology class, which the
closure theorem rules out.

\textbf{8. The Koide Formula}

\textbf{Theorem {[}LV{]}.} \emph{The BDG integers yield the Koide lepton
mass ratio K = 2/3.}

The Koide formula {[}11{]} relates the three charged lepton masses: K =
(m\textsubscript{e} + m\textsubscript{μ} +
m\textsubscript{τ})/((√m\textsubscript{e} + √m\textsubscript{μ} +
√m\textsubscript{τ})²) = 2/3 to high experimental precision (0.01\%). In
the BDG framework, this ratio follows from the SU(3)\textsubscript{gen}
generation structure combined with the BDG depth weighting. The result K
= 2/3 is verified by exact arithmetic in Lean 4 (koide\_formula in
RA\_D1\_Proofs.lean).

\textbf{9. Colour Confinement and Hadron Physics}

\textbf{9.1 Confinement depths {[}LV{]}}

The BDG topology determines the confinement depths of colour-charged
particles: L = 3 for gluons and L = 4 for quarks (confinement\_lengths
in RA\_D1\_Proofs.lean). These are the minimum BDG depths at which
colour-neutral bound states can form.

\textbf{9.2 Proton mass {[}CV{]}}

The proton mass is derived from the BDG Regge structure:

\begin{quote}
\emph{m}
\end{quote}

compared to the experimental value 938.3 MeV (0.08\% agreement). Here
Λ\textsubscript{QCD} = 332 MeV is the standard QCD scale parameter, and
√c\textsubscript{4} = √8 is the BDG vertex weight.

\textbf{9.3 Roper resonance {[}CV{]}}

The Roper resonance N(1440) has a 290 MeV gap above the bare BDG
prediction of 1150 MeV. This gap is identified as a two-pion loop (a
Kind 2 virtual process in BDG language): the bare BDG mass of 1150 MeV
receives a 290 MeV self-energy correction from the ππ intermediate
state, giving the physical mass 1440 MeV.

\textbf{10. The Baryon-to-Dark-Matter Density Ratio}

The cosmological density ratio Ω\textsubscript{CDM}/Ω\textsubscript{b}
is derived from BDG path weights at μ = 1:

\begin{quote}
\emph{f₀ = W}
\end{quote}

where W\textsubscript{baryon} = 3/2 exactly and
W\textsubscript{other}/W\textsubscript{baryon} = 17.32 are proved from
BDG combinatorics at μ = 1 {[}CV{]}. The Planck 2018 measurement is
5.416 ± 0.015, giving 0.3\% agreement. This is the first derivation of a
cosmological density ratio from pure causal-graph topology.

\textbf{11. Gauge Group Structure}

The Standard Model gauge group SU(3) × SU(2) × U(1) emerges from the
alternating sign structure of the BDG coefficients.

\textbf{11.1 The sign mechanism (GS02)}

The BDG coefficients have signs (+, −, +, −, +), forced by the
inclusion-exclusion counting of causal intervals. At each depth k, the
sign determines whether the corresponding layer contributes
isotropically (+) or anisotropically (−) to the local causal geometry.

The alternating isotropy/anisotropy pattern at successive depths
produces: an SU(3) colour structure from the k = 3, 4 layers (three-fold
anisotropy), an SU(2) weak isospin structure from the k = 1, 2 layers
(two-fold anisotropy), and a U(1) hypercharge from the k = 0 layer
(single isotropic factor). The product SU(3) × SU(2) × U(1) is verified
by exact enumeration of the BDG sign pattern {[}CV{]}.

\textbf{12. Discussion}

\textbf{12.1 Numerology vs. derivation}

The results of this paper differ from parameter-fitting or numerological
coincidences in three structural ways. First, the BDG integers are not
fitted to data --- they are determined by the unique self-consistent
growth rule (Section 3). Second, the same integers simultaneously
determine multiple independent physical quantities (two coupling
constants, the particle spectrum, the mass formula, the cosmological
ratio) with no adjustable parameters. Third, the results are
falsifiable: any single quantity deviating from its BDG-predicted value
would invalidate the entire framework.

\textbf{12.2 Relation to the causal set programme}

This work extends the Benincasa-Dowker-Glaser programme {[}1,2,3{]} by
(i) providing an intrinsic characterisation of the BDG coefficients via
the Yeats chord diagram interpretation and Möbius inversion, without
reference to the continuum d'Alembertian; (ii) deriving physical
observables (coupling constants, particle spectrum) directly from the
BDG integers; and (iii) establishing the uniqueness of d = 4 via the
selectivity criterion D4U02.

\textbf{12.3 Status of the Poisson approximation}

The Poisson-CSG model used in Sections 2.3 and 4 is an approximation:
the BDG acceptance filter S \textgreater{} 0 biases the
N\textsubscript{k} distribution away from Poisson (N\textsubscript{2}
inflated by 65\%, N\textsubscript{3} deflated by 85\%). However, the
specific combinations of N\textsubscript{k} that yield physical
observables are insensitive to this bias at the precision achieved
(0.13\% for α\textsubscript{s}, 0.3\% for f\textsubscript{0}). The
uniqueness theorem (Section 3) does not depend on the Poisson model at
all.

\textbf{12.4 Predictions}

The framework makes specific predictions testable with existing or
near-future experiments:

(i) No new stable particle types beyond the Standard Model. The
five-topology closure is machine-verified.

(ii) No proton decay. Baryon number is exactly conserved as the LLC on a
topological winding number.

(iii) θ\textsubscript{QCD} = 0 exactly. The DAG is acyclic; instantons
(requiring Euclidean closed loops) cannot exist.

(iv) Three and only three generations. A fourth generation requires a
sixth topology type, ruled out by L11.

(v) Neutrino mass sum Σmν ≈ 59 meV (from SU(3)\textsubscript{gen} Koide
fit with normal hierarchy). Testable with CMB-S4 and Euclid.

\textbf{13. Conclusions}

We have derived the electromagnetic fine structure constant
α\textsubscript{EM} = 1/137 and the strong coupling constant
α\textsubscript{s}(m\textsubscript{Z}) = 1/√72 from the combinatorial
structure of four-dimensional causal sets, with no free parameters. The
BDG integers (1, −1, 9, −16, 8) that govern these derivations are the
unique self-consistent growth rule for an irreversibly growing directed
acyclic graph, determined by Yeats's chord diagram moment sequence and
Möbius inversion on the Boolean lattice.

From the same integers, we derive the complete Standard Model particle
spectrum (five topology types, machine-verified), the Koide lepton mass
formula K = 2/3, colour confinement depths, the proton mass to 0.08\%,
the Roper resonance mass gap, the gauge group SU(3) × SU(2) × U(1), and
the baryon-to-dark-matter density ratio f\textsubscript{0} = 5.42
(Planck: 5.416). No parameter in any of these derivations is
independently adjustable.

The uniqueness theorem establishes that these results are not one set of
predictions among many --- they are the \emph{only} predictions a
self-consistent causal graph can make. The laws of physics encoded in
the BDG action are not chosen from alternatives. They are the unique
solution to a combinatorial identity.

Companion papers derive further consequences of the same framework: wave
function collapse as irreversible entropy production (Paper II
{[}12{]}), cosmological parameters including the structural zero
cosmological constant and nucleation cosmology (Paper III {[}13{]}), the
Causal Firewall and origin-of-life constraints (Paper IV {[}14{]}), and
the complete framework synthesis (Paper V {[}15{]}). An accessible
narrative overview is given in Paper 0 {[}16{]}.

\textbf{References}

{[}1{]} D. M. T. Benincasa and F. Dowker, Phys. Rev. Lett. 104, 181301
(2010).

{[}2{]} F. Dowker and L. Glaser, Class. Quantum Grav. 30, 195016 (2013).

{[}3{]} L. Glaser, Class. Quantum Grav. 31, 095007 (2014).

{[}4{]} L. Bombelli, J. Lee, D. Meyer, and R. D. Sorkin, Phys. Rev.
Lett. 59, 521 (1987).

{[}5{]} R. D. Sorkin, in Lectures on Quantum Gravity (Springer, 2005).

{[}6{]} D. Rideout and R. D. Sorkin, Phys. Rev. D 61, 024002 (2000).

{[}7{]} J. F. Sandeman, RA\_AmpLocality.lean (2026). GitHub:
github.com/jsandeman/Relational-Actualism.

{[}8{]} G.-C. Rota, Z. Wahrscheinlichkeitstheorie 2, 340--368 (1964).

{[}9{]} K. Yeats, Class. Quantum Grav. 42, 145003 (2025).
arXiv:2412.14036.

{[}10{]} J. F. Sandeman, D4U02 proof document (2026). Zenodo:
doi:10.5281/zenodo.19362289.

{[}11{]} Y. Koide, Lett. Nuovo Cim. 34, 201--205 (1982).

{[}12{]} J. F. Sandeman, Paper II of this series (2026). Zenodo:
doi:10.5281/zenodo.19362968.

{[}13{]} J. F. Sandeman, Paper III of this series (2026). Zenodo:
doi:10.5281/zenodo.19363017.

{[}14{]} J. F. Sandeman, Paper IV of this series (2026).

{[}15{]} J. F. Sandeman, Paper V of this series (2026). Zenodo:
doi:10.5281/zenodo.19363077.

{[}16{]} J. F. Sandeman, Paper 0: A Biography of the Universe (2026).

\end{document}
